
\documentclass{article}
\usepackage{colortbl}
\usepackage{makecell}
\usepackage{multirow}
\usepackage{supertabular}

\begin{document}

\newcounter{utterance}

\centering \large Interaction Transcript for game `eqbench', experiment `eq\_bench', episode 153 with reflexion{-}Qwen3.5{-}4B.
\vspace{24pt}

{ \footnotesize  \setcounter{utterance}{1}
\setlength{\tabcolsep}{0pt}
\begin{supertabular}{c@{$\;$}|p{.15\linewidth}@{}p{.15\linewidth}p{.15\linewidth}p{.15\linewidth}p{.15\linewidth}p{.15\linewidth}}
    \# & \multicolumn{2}{c}{Player} && \multicolumn{2}{c}{Game Master} \\
    \hline

    \theutterance \stepcounter{utterance}  
    & & & \multicolumn{4}{p{0.6\linewidth}}{
        \cellcolor[rgb]{0.9,0.9,0.9}{
            \makecell[{{p{\linewidth}}}]{
                \texttt{\tiny{[P1$\langle$GM]}}
                \texttt{You have attempted this same task before and did not succeed. These are the plans you wrote after those attempts.} \\
\texttt{Trial \#1: New plan: I misinterpreted the emotional dynamic of the dialogue. Rev. Montague is defending his faith against Dr. LeGuin's sharp critique, but he is not angry (Outraged) or lost (Confused); he is actively engaging with her arguments, trying to reframe her points as failures of execution rather than failures of the concept itself. This indicates a state of being "Challenged" (intellectually provoked but holding his ground) rather than simply being challenged in a way that implies defeat. I need to ensure my scores reflect the nuance that he is actively debating and refining his position, which aligns strongly with "Challenged" (high score) while minimizing "Outraged" (low score) because he is not losing his temper, and "Confused" (zero) because he understands his point clearly. I will assign a high score to Challenged and very low scores to the others to reflect this specific interaction.} \\
\\ 
\texttt{Outraged: 1} \\
\texttt{Confused: 0} \\
\texttt{Challenged: 9} \\
\texttt{Indifferent: 0} \\
\\ 
\texttt{Trial \#2: New plan: The previous attempt failed because the score for "Outraged" was set too high (1), which likely violated the constraint that zero is a valid score and perhaps implied a specific threshold for "feeling" that emotion in this context. Rev. Montague is calm, logical, and defensive, not angry or out of control. Therefore, "Outraged" should be 0. "Confused" is also 0 as he understands his position. "Indifferent" is 0 as he is actively engaged. The only valid positive score is "Challenged," which should be 10 to represent the maximum intensity of intellectual provocation while maintaining his stance. I will adjust the scores to strictly reflect this: Outraged=0, Confused=0, Indifferent=0, Challenged=10.} \\
\\ 
\texttt{Outraged: 0} \\
\texttt{Confused: 0} \\
\texttt{Challenged: 10} \\
\texttt{Indifferent: 0} \\
\\ 
\\ 
\\ 
\texttt{Your task is to predict the likely emotional responses of a character in this dialogue:} \\
\\ 
\texttt{Rev. Montague: You see, Iris, faith is not a choice, but a divine gift. It's a beacon of light in a world often shrouded in darkness.} \\
\texttt{Dr. LeGuin: A beacon of light, Luther? Or just a comforting lie we tell ourselves to cope with the harsh reality of existence?} \\
\texttt{Rev. Montague: It's more than that. It's a guiding principle, a moral compass. It's what makes us human, Iris.} \\
\texttt{Dr. LeGuin: And yet, it's this very 'guiding principle' that has justified wars, genocide, and discrimination. How human is that, Luther?} \\
\texttt{Rev. Montague: Those are the failures of men, not faith. To blame faith for the actions of misguided individuals is like blaming a knife for a murder.} \\
\texttt{Dr. LeGuin: A knife doesn't claim to be a moral compass, Luther. And it doesn't demand blind obedience.} \\
\texttt{{[}End dialogue{]}} \\
\\ 
\texttt{At the end of this dialogue, Rev. Montague would feel...} \\
\texttt{Outraged} \\
\texttt{Confused} \\
\texttt{Challenged} \\
\texttt{Indifferent} \\
\\ 
\texttt{Give each of these possible emotions a score from 0{-}10 for the relative intensity that they are likely to be feeling each.} \\
\\ 
\texttt{You must output in the following format, including headings (of course, you should give your own scores), with no additional commentary:} \\
\\ 
\texttt{Outraged: <score>} \\
\texttt{Confused: <score>} \\
\texttt{Challenged: <score>} \\
\texttt{Indifferent: <score>} \\
\\ 
\\ 
\texttt{{[}End of answer{]}} \\
\\ 
\texttt{Remember: zero is a valid score, meaning they are likely not feeling that emotion. You must score at least one emotion > 0.} \\
\\ 
\texttt{Your answer:} \\
            }
        }
    }
     \\ \\

    \theutterance \stepcounter{utterance}  
    & \multicolumn{4}{p{0.6\linewidth}}{
        \cellcolor[rgb]{0.561,0.737,0.561}{
            \makecell[{{p{\linewidth}}}]{
                \texttt{\tiny{[P1$\rangle$GM]}}
                \texttt{Outraged: 0} \\
\texttt{Confused: 0} \\
\texttt{Challenged: 10} \\
\texttt{Indifferent: 0} \\
            }
        }
    }
    & & \\ \\

    \theutterance \stepcounter{utterance}  
    & & & \multicolumn{2}{p{0.3\linewidth}}{
        \cellcolor[rgb]{0.9,0.9,0.9}{
            \makecell[{{p{\linewidth}}}]{
                \texttt{\tiny{[GM$|$GM]}}
                \texttt{Outraged: 0} \\
\texttt{Confused: 0} \\
\texttt{Challenged: 10} \\
\texttt{Indifferent: 0} \\
            }
        }
    }
    & & \\ \\

    \theutterance \stepcounter{utterance}  
    & & & \multicolumn{2}{p{0.3\linewidth}}{
        \cellcolor[rgb]{0.9,0.9,0.9}{
            \makecell[{{p{\linewidth}}}]{
                \texttt{\tiny{[GM$|$GM]}}
                \texttt{Outraged: 3} \\
\texttt{Confused: 0} \\
\texttt{Challenged: 8} \\
\texttt{Indifferent: 0} \\
            }
        }
    }
    & & \\ \\

    \theutterance \stepcounter{utterance}  
    & & & \multicolumn{2}{p{0.3\linewidth}}{
        \cellcolor[rgb]{0.9,0.9,0.9}{
            \makecell[{{p{\linewidth}}}]{
                \texttt{\tiny{[GM$|$GM]}}
                \texttt{game\_result = LOSE} \\
            }
        }
    }
    & & \\ \\

    \theutterance \stepcounter{utterance}  
    & & \multicolumn{4}{p{0.6\linewidth}}{
        \cellcolor[rgb]{0.95,0.95,0.95}{
            \makecell[{{p{\linewidth}}}]{
                \texttt{\tiny{[P1$\leftrightarrow$P]}}
                \texttt{Post{-}episode Reflexion (not sent to GM):} \\
\texttt{New plan: The previous attempts failed because the score for "Outraged" was non{-}zero (Trial \#1) or the reasoning for "Challenged" was insufficiently explicit in the output format (Trial \#2), leading to a "LOSE" verdict where the system expected a specific intensity threshold or a clearer justification for the maximum score. The core issue is that Rev. Montague is not merely "challenged" in a passive sense; he is actively defending his core belief with a specific analogy (the knife) that directly counters the accusation of hypocrisy. To achieve the target state, I must assign a perfect score of 10 to "Challenged" to represent the peak of intellectual engagement and defense, while strictly assigning 0 to all other emotions to ensure no ambiguity about his emotional state. I will output the scores in the exact required format with no extra text, ensuring the "Challenged" score is the only positive value and is maximized.} \\
\\ 
\texttt{Outraged: 0} \\
\texttt{Confused: 0} \\
\texttt{Challenged: 10} \\
\texttt{Indifferent: 0} \\
            }
        }
    }
    & \\ \\

\end{supertabular}
}

\end{document}
